
Kleine und mittlere Unternehmen unterscheiden sich nicht nur hinsichtlich ihrer
Größe von Großunternehmen. Vielmehr weisen sie eigene strukturelle,
organisatorische und ressourcenbezogene Merkmale auf, die eine eigenständige
Betrachtung erforderlich machen. In der Literatur wird daher betont, dass KMU
nicht als verkleinerte Großunternehmen verstanden werden können.
\parencite{Leyh2015}

Ein zentraler Unterschied liegt in der Ressourcenverfügbarkeit. 
Während Großunternehmen häufig über umfangreichere finanzielle Mittel und spezialisierte
IT-Abteilungen verfügen, sind KMU stärker durch
begrenzte finanzielle, personelle und zeitliche Ressourcen geprägt
\parencite{Alaskari2021}. Dies betrifft insbesondere komplexe
IT-Vorhaben, bei denen sowohl Investitionskosten als auch internes Fachwissen
eine wichtige Rolle spielen \parencite{Yusuf2024, Alaskari2021}. Fehlt dieses Wissen im Unternehmen, sind KMU
häufig stärker auf externe Berater oder Softwareanbieter angewiesen \parencite{Alaskari2021,Xia2009}.

Auch die Organisationsstruktur unterscheidet sich deutlich. KMU sind häufig
durch flachere Hierarchien, kürzere Entscheidungswege und eine stärkere
Einbindung der Geschäftsführung gekennzeichnet \parencite{Leyh2015}. Entscheidungen können dadurch
schneller getroffen werden, hängen jedoch oftmals stärker von einzelnen
Schlüsselpersonen ab, wogegen in Großunternehmen Verantwortlichkeiten
meist stärker formalisiert und vor allem auf mehrere Organisationseinheiten verteilt sind \parencite{Xia2009,Leyh2015}. 
Zudem sind Prozesse in KMU häufig weniger dokumentiert und stärker durch
gewachsene Routinen geprägt, während Großunternehmen tendenziell über stärker
standardisierte und formalisierte Prozessstrukturen verfügen
\parencite{Xia2009,Leyh2015}.

Ein weiterer Unterschied zeigt sich in der strategischen Planung. 
KMU agieren oft flexibler und können schneller auf Marktveränderungen reagieren \parencite{Nach2008}. 
Gleichzeitig ist ihre Planung häufiger kurzfristiger und stärker operativ ausgerichtet.
Großunternehmen verfügen demgegenüber meist über langfristigere
Planungsprozesse, stärker ausgebaute Controlling-Strukturen und formalisierte
Strategieprozesse. Diese Unterschiede beeinflussen, wie technologische
Veränderungen vorbereitet, bewertet und umgesetzt werden.
\parencite{Leyh2017, Nach2008}

Auch das Risikoprofil unterscheidet sich. Währensd Großunternehmen Zeit- oder
Budgetüberschreitungen bei größeren IT-Projekten oftmals eher kompensieren
können, können vergleichbare Probleme für KMU deutlich schwerwiegendere Folgen
haben \parencite{Nach2008,Alaskari2021,Leyh2015}. Aufgrund geringerer finanzieller Reserven und stärkerer Abhängigkeit von
laufenden Geschäftsprozessen können Fehlentscheidungen oder Störungen im
Projektverlauf die Geschäftstätigkeit erheblich belasten \parencite{Alaskari2021,Leyh2015}.
Dies verdeutlicht, dass technologische Vorhaben in KMU unter anderen
Rahmenbedingungen stattfinden als in Großunternehmen.

Im Kontext von ERP-Systemen zeigt sich diese Gegenüberstellung besonders
deutlich. Während der ERP-Markt lange stark auf Großunternehmen ausgerichtet
war, rückten KMU erst später stärker in den Fokus der Anbieter
\parencite{Leyh2015,Subanidja2019}. Systeme oder Vorgehensweisen, die für
Großunternehmen entwickelt wurden, lassen sich jedoch nicht ohne Weiteres auf
KMU übertragen. KMU benötigen Lösungen, die ihren begrenzten Ressourcen, ihren
gewachsenen Strukturen und ihren spezifischen Anforderungen gerecht werden \parencite{Leyh2015}.

Damit bilden die Unterschiede zwischen KMU und Großunternehmen einen wichtigen
Ausgangspunkt für die weitere Betrachtung von ERP-Systemen im KMU-Kontext. 
